% GENERATED FILE. Edit canonical lessons, then run npm run generate:books.
% canonical-source-hash: fnv1a64:80f7f20bf4014410
% canonical-lessons: GU-C216-tapvu, GU-C216-ukalat, GU-C216-bhej, GU-C216-tapman, GU-C216-digri

\chapter{The hot months}
\label{ch:gu-the-hot-months}
\hlchaptermodality{216}

\begin{chapteropening}
\textbf{By the end of this chapter:} \emph{I can say how hot it will get and what the temperature will be.}

Complete the last lesson of chapter 216: I can say how hot it will get and what the temperature will be.
\end{chapteropening}

\section[\texorpdfstring{\gu{તપવું} (tapvũ)}{તપવું (tapvũ)}]{\gu{તપવું} (tapvũ) — to blaze, to be hot}

\label{lesson:GU-C216-tapvu}



\begin{quote}
Before the new one: say the Gujarati for thick, then the Gujarati for lukewarm.
\end{quote}

\subsection*{You'll want to know: \gu{તપવું}}
\textbf{\gu{તપવું}} — \emph{tapvũ} — \textquotedblleft{}to blaze, to be hot\textquotedblright{}.

\begin{quote}
\textbf{\gu{બપોરે} \gu{સૂરજ} \gu{તપશે}.} — \emph{bapore sūraj tapśe} — \textquotedblleft{}The sun will blaze at midday.\textquotedblright{}
\end{quote}

\begin{cousinweb}
\textbf{\gu{તપવું}} (\emph{tapvũ}) is usually traced to Sanskrit \emph{tapyati}, \textquotedblleft{}is hot\textquotedblright{}. CDIAL does not list \emph{tapvũ} itself, but it lists the related Gujarati \emph{tāpvũ}, \textquotedblleft{}to warm oneself\textquotedblright{}, under the same entry (CDIAL 5684).
\end{cousinweb}

\begin{grammarlens}[title={What you've built}]
One more word for everyday life at A2.
\end{grammarlens}

\subsection*{Your turn}
\begin{itemize}
\raggedright
  \item \emph{Say it:} \emph{tapvũ}
  \item \emph{Say it:} \emph{tapvũ}, once more
  \item \emph{Say it:} \emph{bapore sūraj tapśe} — The sun will blaze at midday.
  \item \emph{Recall:} say the Gujarati for thick, then the Gujarati for lukewarm, then say \emph{tapvũ} again
\end{itemize}

\begin{tcolorbox}[breakable,colback=teal!4,colframe=teal!35!black,title={Before you move on}]
What does \gu{તપવું} mean? (\textquotedblleft{}To blaze, to be hot\textquotedblright{}.) Say it once more.
\end{tcolorbox}
\section[\texorpdfstring{\gu{ઉકળાટ} (ukaḷāṭ)}{ઉકળાટ (ukaḷāṭ)}]{\gu{ઉકળાટ} (ukaḷāṭ) — sultry heat, mugginess}

\label{lesson:GU-C216-ukalat}



\begin{quote}
Before the new one: say the Gujarati for lukewarm, then the Gujarati for to blaze.
\end{quote}

\subsection*{You'll want to know: \gu{ઉકળાટ}}
\textbf{\gu{ઉકળાટ}} — \emph{ukaḷāṭ} — \textquotedblleft{}sultry heat, mugginess\textquotedblright{}.

\begin{quote}
\textbf{\gu{બપોરે} \gu{ઘણો} \gu{ઉકળાટ} \gu{હશે}.} — \emph{bapore ghaṇo ukaḷāṭ haśe} — \textquotedblleft{}It will be very muggy at midday.\textquotedblright{}
\end{quote}

\begin{cousinweb}
\textbf{\gu{ઉકળાટ}} (\emph{ukaḷāṭ}) is built on the verb \emph{ukaḷvũ}, \textquotedblleft{}to boil\textquotedblright{}, so muggy air is air that seems to boil. CDIAL derives \emph{ukaḷvũ} from \emph{*utkalati}, \textquotedblleft{}goes out or up\textquotedblright{}, a reconstructed form (CDIAL 1716).
\end{cousinweb}

\begin{grammarlens}[title={What you've built}]
One more word for everyday life at A2.
\end{grammarlens}

\subsection*{Your turn}
\begin{itemize}
\raggedright
  \item \emph{Say it:} \emph{ukaḷāṭ}
  \item \emph{Say it:} \emph{ukaḷāṭ}, once more
  \item \emph{Say it:} \emph{bapore ghaṇo ukaḷāṭ haśe} — It will be very muggy at midday.
  \item \emph{Recall:} say the Gujarati for lukewarm, then the Gujarati for to blaze, then say \emph{ukaḷāṭ} again
\end{itemize}

\begin{tcolorbox}[breakable,colback=teal!4,colframe=teal!35!black,title={Before you move on}]
What does \gu{ઉકળાટ} mean? (\textquotedblleft{}Sultry heat, mugginess\textquotedblright{}.) Say it once more.
\end{tcolorbox}
\section[\texorpdfstring{\gu{ભેજ} (bhej)}{ભેજ (bhej)}]{\gu{ભેજ} (bhej) — humidity, damp}

\label{lesson:GU-C216-bhej}



\begin{quote}
Before the new one: say the Gujarati for to blaze, then the Gujarati for sultry heat.
\end{quote}

\subsection*{You'll want to know: \gu{ભેજ}}
\textbf{\gu{ભેજ}} — \emph{bhej} — \textquotedblleft{}humidity, damp\textquotedblright{}.

\begin{quote}
\textbf{\gu{ચોમાસામાં} \gu{ઘણો} \gu{ભેજ} \gu{હશે}.} — \emph{comāsāmā̃ ghaṇo bhej haśe} — \textquotedblleft{}There will be a lot of humidity in the monsoon.\textquotedblright{}
\end{quote}

\begin{cousinweb}
The origin of \textbf{\gu{ભેજ}} (\emph{bhej}) is not settled. CDIAL lists no Gujarati form for it, and none of its entries for damp or moisture holds the word.
\end{cousinweb}

\begin{grammarlens}[title={What you've built}]
One more word for everyday life at A2.
\end{grammarlens}

\subsection*{Your turn}
\begin{itemize}
\raggedright
  \item \emph{Say it:} \emph{bhej}
  \item \emph{Say it:} \emph{bhej}, once more
  \item \emph{Say it:} \emph{comāsāmā̃ ghaṇo bhej haśe} — There will be a lot of humidity in the monsoon.
  \item \emph{Recall:} say the Gujarati for to blaze, then the Gujarati for sultry heat, then say \emph{bhej} again
\end{itemize}

\begin{tcolorbox}[breakable,colback=teal!4,colframe=teal!35!black,title={Before you move on}]
What does \gu{ભેજ} mean? (\textquotedblleft{}Humidity, damp\textquotedblright{}.) Say it once more.
\end{tcolorbox}
\section[\texorpdfstring{\gu{તાપમાન} (tāpmān)}{તાપમાન (tāpmān)}]{\gu{તાપમાન} (tāpmān) — temperature}

\label{lesson:GU-C216-tapman}



\begin{quote}
Before the new one: say the Gujarati for sultry heat, then the Gujarati for humidity.
\end{quote}

\subsection*{You'll want to know: \gu{તાપમાન}}
\textbf{\gu{તાપમાન}} — \emph{tāpmān} — \textquotedblleft{}temperature\textquotedblright{}.

\begin{quote}
\textbf{\gu{શિયાળામાં} \gu{તાપમાન} \gu{નીચું} \gu{હશે}.} — \emph{śiyāḷāmā̃ tāpmān nīcũ haśe} — \textquotedblleft{}The temperature will be low in winter.\textquotedblright{}
\end{quote}

\begin{cousinweb}
\textbf{\gu{તાપમાન}} (\emph{tāpmān}) is a learned word built from Sanskrit \emph{tāpa}, \textquotedblleft{}heat\textquotedblright{} (CDIAL 5767), and \emph{māna}, \textquotedblleft{}measure\textquotedblright{}; EtymDB lists it as borrowed from Sanskrit.
\end{cousinweb}

\begin{grammarlens}[title={What you've built}]
One more word for everyday life at A2.
\end{grammarlens}

\subsection*{Your turn}
\begin{itemize}
\raggedright
  \item \emph{Say it:} \emph{tāpmān}
  \item \emph{Say it:} \emph{tāpmān}, once more
  \item \emph{Say it:} \emph{śiyāḷāmā̃ tāpmān nīcũ haśe} — The temperature will be low in winter.
  \item \emph{Recall:} say the Gujarati for sultry heat, then the Gujarati for humidity, then say \emph{tāpmān} again
\end{itemize}

\begin{tcolorbox}[breakable,colback=teal!4,colframe=teal!35!black,title={Before you move on}]
What does \gu{તાપમાન} mean? (\textquotedblleft{}Temperature\textquotedblright{}.) Say it once more.
\end{tcolorbox}
\section[\texorpdfstring{\gu{ડિગ્રી} (ḍigrī)}{ડિગ્રી (ḍigrī)}]{\gu{ડિગ્રી} (ḍigrī) — degree}

\label{lesson:GU-C216-digri}



\begin{quote}
Before the new one: say the Gujarati for humidity, then the Gujarati for temperature.
\end{quote}

\subsection*{You'll want to know: \gu{ડિગ્રી}}
\textbf{\gu{ડિગ્રી}} — \emph{ḍigrī} — \textquotedblleft{}degree\textquotedblright{}.

\begin{quote}
\textbf{\gu{કાલે} \gu{ચાળીસ} \gu{ડિગ્રી} \gu{હશે}.} — \emph{kāle cāḷīs ḍigrī haśe} — \textquotedblleft{}It will be forty degrees tomorrow.\textquotedblright{}
\end{quote}

A feminine noun.

\begin{cousinweb}
\textbf{\gu{ડિગ્રી}} (\emph{ḍigrī}) is a loan from English \emph{degree}, for heat and for a university degree alike.
\end{cousinweb}

\begin{grammarlens}[title={What you've built}]
One more word for everyday life at A2.
\end{grammarlens}

\subsection*{Your turn}
\begin{itemize}
\raggedright
  \item \emph{Say it:} \emph{ḍigrī}
  \item \emph{Say it:} \emph{ḍigrī}, once more
  \item \emph{Say it:} \emph{kāle cāḷīs ḍigrī haśe} — It will be forty degrees tomorrow.
  \item \emph{Recall:} say the Gujarati for humidity, then the Gujarati for temperature, then say \emph{ḍigrī} again
\end{itemize}

\begin{tcolorbox}[breakable,colback=teal!4,colframe=teal!35!black,title={Before you move on}]
What does \gu{ડિગ્રી} mean? (\textquotedblleft{}Degree\textquotedblright{}.) Say it once more.
\end{tcolorbox}
